The results hold for the plants, the sample sizes and the parameters of the experiments, and the threats to their validity are of three kinds.

\subsection{Internal validity}

Every rate is a Monte Carlo estimate with a standard error of at most 0.025 (the resilience and situation rates near 0.5) and of about 0.003 to 0.008 for the coverage rates, so differences smaller than twice the standard error of the difference are not claimed. Several predictions are judged by comparing a rate with the level plus twice its standard error, a bound that has a false-rejection probability of about one in forty on each comparison; the manuscript reports the ranges over grids and names the cells where a difference is within the error, but it applies no correction for the number of comparisons. The coverage experiment evaluates the variants on the covered count, which is the only statistic that the tests read; the vectorised p-value tables are checked against the engine function on every replication (\CovEngineMismatches{} mismatches over \CovTotalReps{}), and against the exact sizes by summation, so the risk of a bug that the checks do not catch is small, though the resilience and situation experiments call the engine functions directly and have no second implementation to compare with. The permutation reference may not be exactly valid in the engine's construction: its pseudo-development set has fewer rows than the real development set, and the pooled sizes exceed the exact sizes of the p-value grid by about 0.01 (one observation, its cause not tested; \cref{sec:results}).

\subsection{External validity}

The generators are plants chosen to isolate one defect each, and a bank's data realise other nulls and other defects. The coverage test is the exception for its null, which is exact for every continuous score distribution; its power is measured for one family of drifts, in which the test scores are raised to a power, and a drift that moves the scores in a different way (a change in the dispersion, a change of the correlation between the score and the label) has a different power. Ties in the scores, which discrete scores produce, are outside \cref{prop:law}: a model with few distinct scores gives a set that covers more rows than the continuous law allows, and the Beta-Binomial test is then not exact. The windows are assumed exchangeable with the calibration rows under the null; a time series of out-of-time windows has autocorrelation, and the test is then oversized by an amount that the experiments do not measure. The resilience experiment fixes the development sample at twice the evaluation sample, the worst share at 0.25 (the committed run uses 0.10), one model and one shift direction. The situation experiments use 30 neighbours, one distance and two designs, and sizes of 400 and 1,000 applicants; the committed fairness evidence has \CmFairQMin{} to \CmFairQMax{} complainants, and the behaviour of the null share between 1,000 applicants and that size is not measured.

\subsection{Construct validity}

A rejection counts as a detection when the plant has a defect, and as a false alarm when it has none, and the defects are the ones the plants plant. The shift of the resilience plant is a shift in the covariates in a region where the true model changes, so that the model fails on the shifted rows by construction; a shift that does not change the model's ranking is not a defect, and the test is not designed to detect it. The bias of the situation plants is a lower cut-off or a higher probability of a negative decision for the protected group; a bias that acts only through attributes that the test matches on is not detected by a ceteris-paribus comparator, and the manuscript does not claim that it is. The size of the situation light is measured against the absence of a group effect, and the light reports no legal conclusion; a green light is not a finding of fairness and an amber light is not a finding of discrimination.

\subsection{Scope}

The manuscript does not claim that the Beta-Binomial law is new; it is published \citep{vovk2012,angelopoulos2021}. It does not claim that the permutation reference is optimal; it compares it with four other rules on one plant. It does not claim the cut 0.20 is the best cut for any portfolio; it shows that the size of a cut has to be measured on a plant that resembles the portfolio. The committed rows of the product are shown as the behaviour of the corrected rules on the product's data and not as a measurement of any error rate, because their truth is not planted. Conformal prediction under covariate shift and under online updating \citep{tibshirani2019,gibbs2021,barber2023} changes the null of the coverage test, and the tests of this manuscript do not apply to a set that is built for a shift. The sizes of the corrected rules are measured in the generators of the engine at the settings of the experiments, and a bank that applies the rules to its own data should repeat the measurement with a plant that has the sample sizes and the structure of that data, using the protocol of \cref{alg:protocol}.
